\section*{Bab \ref{ch:b1:clastic-sediments} --- Sedimen klastik dan lingkungannya}

\begin{solution}{exo:b1:clastic-sediments:1}
$\varphi = -\log_2 d$: \qty{0.18}{mm} memberi $2.47$, pasir halus (2 hingga 3);
\qty{0.75}{mm} memberi $0.42$, pasir kasar (0 hingga 1); \qty{3}{mm} memberi
$-1.58$, granul ($-1$ hingga $-2$); \qty{0.012}{mm} memberi $6.38$, lanau halus (6
hingga 7); \qty{150}{mm} memberi $-7.23$, kerakal ($-6$ hingga $-8$).
\end{solution}

\begin{solution}{exo:b1:clastic-sediments:2}
$M_z = (1.05 + 1.71 + 2.43)/3 = 1.73\,\varphi$, yakni $2^{-1.73} =
\qty{0.30}{mm}$, pasir sedang. $\sigma_I = (2.43 - 1.05)/4 + (2.98 - 0.62)/6.6
= 0.345 + 0.358 = 0.70\,\varphi$: cukup baik terpilah (0.50 hingga 0.71), tepat
di bawah batas agak terpilah.
\end{solution}

\begin{solution}{exo:b1:clastic-sediments:3}
$q_b = (1 - 0.4) \times 0.02 \times (\num{0.01}/60)/2 = \qty{1.0e-6}{m^2.s^{-1}}$
butiran padat. Dalam sehari, $\num{1.0e-6} \times 86400 = \qty{0.0864}{m^3}$ per
meter puncak, atau $0.0864 \times 2650 = \qty{229}{kg}$ kuarsa.
\end{solution}

\begin{solution}{exo:b1:clastic-sediments:4}
(a)~50 knot adalah $50 \times 1852/3600 = \qty{25.7}{m.s^{-1}}$; 12 knot,
\qty{6.17}{m.s^{-1}}. (b)~$\qty{100000}{m}/\qty{6.17}{m.s^{-1}} =
\qty{16200}{s}$, sekitar \qty{4.5}{h}. (c)~480 mil adalah \qty{772}{km}, 400 mil
\qty{644}{km}.
\end{solution}

\begin{solution}{exo:b1:clastic-sediments:5}
(a)~$(\varphi_{84} - \varphi_{16})/2 = 2 \times 0.9945\,\sigma/2 =
0.9945\,\sigma$, dan $(\varphi_{95} - \varphi_5)/3.3 = 2 \times
1.6449\,\sigma/3.3 = 0.9969\,\sigma$: keduanya memperkirakan $\sigma$ dalam sepertiga persen. (b)~Sedimen sungguhan tidak perlu normal dalam satuan phi: campuran kecil butiran yang jauh lebih kasar atau jauh lebih halus mengubah $\varphi_5$ atau $\varphi_{95}$ banyak dan $\varphi_{16}$ serta $\varphi_{84}$ sedikit, sehingga perkiraan kedua tumbuh sementara yang pertama tidak. Keduanya berbeda paling banyak untuk campuran, seperti pasir dengan sedikit kerikil dan lumpur, atau endapan bimodal.
\end{solution}

\begin{solution}{exo:b1:clastic-sediments:6}
(a)~$\sin\theta = T/L = 1.5/15 = 0.1$, sehingga $\theta = \ang{5.74}$. (b)~$a =
c\tan\theta = 0.5 \times 0.1005 = \qty{0.050}{cm}$ per menit, atau
\qty{3.0}{cm} per jam. (c)~Dasar yang dinaikkan beberapa sentimeter per jam menerima pasir yang jatuh dari suspensi dengan sangat cepat: arus bermuatan berat yang sedang melambat, seperti banjir yang surut atau ekor \omterm{def:b1:clastic-sediments:turbidite}{arus turbidit}.
\end{solution}

\begin{solution}{exo:b1:clastic-sediments:7}
(a)~$(350 + 20 + 5 + 340 + 15 + 30)/6 = 760/6 = 126.7$: arah tenggara,
meskipun setiap bacaan berada dalam \ang{20} dari utara. (b)~$S =
\sum\sin\alpha_i = 0.672$, $C = \sum\cos\alpha_i = 5.692$, sehingga $\bar\alpha =
\operatorname{atan2}(0.672, 5.692) = \ang{6.7}$, dan $\bar R =
\sqrt{0.672^2 + 5.692^2}/6 = 5.732/6 = 0.955$. (c)~Rerata aritmetik memperlakukan
350 dan 005 sebagai berjauhan, sedangkan pada lingkaran jaraknya \ang{15}: \omterm{def:b1:spherical-astronomy:horizontal}{azimut}
harus dirata-ratakan sebagai vektor satuan. $\bar R$ yang besar menunjukkan bahwa
bacaan-bacaan itu berkelompok rapat di sekitar \ang{7}.
\end{solution}

\begin{solution}{exo:b1:clastic-sediments:8}
A memiliki lebih dari \qty{5}{\percent} lempung: belum matang, dan secara komposisi juga belum matang, dengan \omterm{def:b1:rock-forming-minerals:groups}{feldspar} dan pecahan batuannya; diendapkan dengan cepat dekat sumbernya, di kipas aluvial atau dataran banjir, atau oleh \omterm{def:b1:clastic-sediments:turbidite}{arus turbidit}. B memiliki kurang dari \qty{5}{\percent} lempung dan terpilah baik ($\sigma_I < 0.5\,\varphi$) tetapi butirannya sub-bersudut ($\rho < 3$): matang, misalnya gosong sungai atau pasir paparan yang diayak oleh arus. C tidak memiliki lempung, terpilah sangat baik dan butirannya membulat ($\rho > 3$): sangat matang, dan dengan \qty{98}{\percent} kuarsa juga matang secara komposisi, seperti di pantai atau di medan \omterm{def:b1:clastic-sediments:bedform}{bukit pasir} gurun.
\end{solution}

\begin{solution}{exo:b1:clastic-sediments:9}
(a)~$1.5/\sqrt{9.81 \times 0.3} = 0.87$; $1.0/\sqrt{9.81 \times 2} = 0.23$.
(b)~$h = U^2/g = 1.5^2/9.81 = \qty{0.23}{m}$; $U = \sqrt{9.81 \times 2} =
\qty{4.43}{m.s^{-1}}$. (c)~Bilangan Froude naik ketika kedalaman turun pada kecepatan tertentu: lembaran air yang tipis dan cepat, seperti limpasan balik (backwash) ombak di pantai, dengan mudah mencapai bilangan Froude mendekati 1, sedangkan sungai yang dalamnya beberapa meter memerlukan beberapa meter per detik, yang hanya dicapai sedikit sungai kecuali saat banjir.
\end{solution}

\begin{solution}{exo:b1:clastic-sediments:10}
(a)~Dengan titik tengah $-0.5, 0.5, \ldots, 4.5$: $m = (-1 + 4 + 45 + 100 + 52.5
+ 22.5)/100 = 2.23\,\varphi$; jumlah selisih kuadrat yang diberi bobot persentase adalah \num{107.71}, sehingga $s = \sqrt{1.0771} = 1.04\,\varphi$:
terpilah buruk. (b)~Persentase kumulatif yang lebih kasar adalah 0, 2, 10, 40, 80, 95 dan 100 pada $-1, 0, \ldots, 5\,\varphi$. Maka $\varphi_5 = 0 + 3/8 = 0.375$,
$\varphi_{16} = 1 + 6/30 = 1.20$, $\varphi_{50} = 2 + 10/40 = 2.25$,
$\varphi_{84} = 3 + 4/15 = 3.27$, $\varphi_{95} = 4.0$. $M_z = (1.20 + 2.25 +
3.27)/3 = 2.24\,\varphi$ dan $\sigma_I = (3.27 - 1.20)/4 + (4.0 - 0.375)/6.6 =
0.517 + 0.549 = 1.07\,\varphi$. (c)~Kedua metode sepakat sekitar
\num{0.01} pada rerata dan \num{0.03} pada \omterm{def:b1:clastic-sediments:sorting}{pemilahan}: pasir halus
($2.24\,\varphi$ adalah \qty{0.21}{mm}), terpilah buruk.
\end{solution}

\begin{solution}{exo:b1:clastic-sediments:11}
(a)~$U = \sqrt{0.02 \times 9.81 \times 50 \times 0.005/0.003} =
\qty{4.04}{m.s^{-1}}$, atau 7.9 knot. (b)~$(\Delta\rho/\rho)\,h = U^2
C_f/(gS) = 25.7^2 \times 0.003/(9.81 \times 0.005) = \qty{40}{m}$: misalnya
arus setebal \qty{400}{m} dengan kelebihan massa jenis \qty{10}{\percent}, yang
menunjukkan betapa besar dan padatnya arus 1929 pastilah. (c)~Kecepatan
berskala sebagai $\sqrt S$: $4.04/\sqrt 5 = \qty{1.81}{m.s^{-1}}$. Arus yang
lebih lambat memiliki turbulensi yang lebih lemah dan menjatuhkan butiran yang lebih kasar; ia menjadi kurang padat, melambat lebih jauh, dan menyebarkan muatannya di atas dataran sebagai \omterm{def:b1:clastic-sediments:turbidite}{lapisan bergradasi}.
\end{solution}

\begin{solution}{exo:b1:clastic-sediments:12}
(a)~$0.7 \times \qty{5e10}{kg}/\qty{2000}{kg.m^{-3}} = \qty{1.75e7}{m^3}$ setahun. (b)~Tersebar pada bagian depan setinggi \qty{30}{m} dan selebar \qty{100}{km}, ia memajukan bagian depan sebesar $\num{1.75e7}/(30 \times \num{1e5}) = \qty{5.8}{m}$ setahun. (c)~$a = v\tan\alpha = 5.83 \times \tan\ang{0.5} = \qty{0.051}{m}$, sekitar \qty{5}{cm} setahun, seratus kali lebih kecil daripada kemajuannya; urutan \qty{30}{m} terbentuk dalam $30/0.051 = 590$ tahun.
\end{solution}

\begin{solution}{pb:b1:clastic-sediments:1}
\textbf{1.} Satuan~1: $(5.0 + 6.6 + 8.4)/3 = 6.67\,\varphi$, \qty{0.0098}{mm}
(\qty{9.8}{\micro m}); satuan~2: $3.47\,\varphi$, \qty{0.090}{mm}; satuan~3:
$1.80\,\varphi$, \qty{0.29}{mm}; satuan~4: $0.73\,\varphi$, \qty{0.60}{mm}.

\textbf{2.} Satuan~1: $(8.4 - 5.0)/4 + (9.6 - 4.2)/6.6 = 0.850 + 0.818 =
1.67\,\varphi$; satuan~2: $0.550 + 0.545 = 1.10$; satuan~3: $0.275 + 0.273 = 0.55$;
satuan~4: $0.725 + 0.712 = 1.44$.

\textbf{3.} Satuan~1: lanau halus, terpilah buruk; satuan~2: pasir sangat halus, terpilah buruk; satuan~3: pasir sedang, cukup baik terpilah; satuan~4: pasir kasar, terpilah buruk.

\textbf{4.} Satuan~3. Di muara sungai arus dan ombak menggerakkan pasir lagi dan lagi serta membawa pergi yang halus, sementara lumpur prodelta mengendap dari suspensi sebagai campuran semua ukuran halus yang sampai padanya.

\textbf{5.} Ukuran $8\,\varphi$ berada pada $1.4/1.8$ jarak antara
$\varphi_{50}$ dan $\varphi_{84}$, sehingga $50 + (1.4/1.8) \times 34 = 76$ persen dari sampel lebih kasar dan sekitar \qty{24}{\percent} berukuran lempung. Partikel halus seperti itu hanya mengendap dari air yang hampir diam: lokasinya berada di bawah jangkauan ombak dan arus.

\textbf{6.} Sebuah saluran membawa kerikil, yang digulirkan di dasarnya saat banjir, sekaligus pasir; endapan itu mencampur lag kerikil dari aliran terdalam dengan pasir dari aliran yang lebih rendah, sehingga ia mencakup banyak kelas.

\textbf{7.} Rerata tumbuh dari \qty{0.0098}{mm} ke \qty{0.090}{mm} dan \qty{0.29}{mm}: suksesi itu mengasar ke atas.

\textbf{8.} Satuan~1, lumpur berlaminasi: prodelta; satuan~2, lanau dan pasir yang berselang-seling: bagian depan delta bawah; satuan~3, pasir berperlapisan silang: gosong muara dari bagian depan delta atas.

\textbf{9.} Setiap zona tiba menempati tempat zona yang dahulu berada di depannya: delta membangun maju, ke arah laut, di atas lokasi itu.

\textbf{10.} Setidaknya \qty{21}{m}, tebal satuan 1 hingga 3, karena suksesi mengisi air dari dasar prodelta hingga gosong muara dekat permukaan; lebih dari itu, karena lumpur telah terpadatkan sejak itu (\cref{ch:b1:diagenesis-basins}).

\textbf{11.} Sebuah saluran distributari dari dataran delta, yang menggerus gosong muara ketika delta membangun maju; kerikilnya adalah \omterm{def:b1:erosion-transport:load}{muatan dasarnya}.

\textbf{12.} Tidak: gosong tepi dimulai dengan permukaan erosi dan lag, dan menghalus ke atas, hingga lumpur dataran banjir. Bagian depan delta mengasar ke atas dari lumpur.

\textbf{13.} $\qty{16}{m}/\num{4000}\text{ tahun} = \qty{4.0}{mm}$ setahun.

\textbf{14.} $v = a/\tan\alpha = 0.0040/\tan\ang{0.5} = 0.0040/0.008727 =
\qty{0.458}{m}$ setahun.

\textbf{15.} $0.458 \times 4000 = \qty{1830}{m}$, sekitar \qty{1.8}{km}.

\textbf{16.} $\qty{21}{m}/\qty{4.0}{mm.yr^{-1}} = 5250$ tahun.

\textbf{17.} Bagian depan delta menerima sebagian besar pasir di dekat muara, prodelta jauh lebih sedikit dan dataran delta sedikit kecuali saat banjir; muara yang aktif juga berpindah-pindah tempat, sehingga suatu titik dapat menerima banyak untuk sementara lalu hampir tidak ada. Satu laju di antara dua lapisan yang bertarikh adalah rerata atas perubahan-perubahan ini, dan lumpur telah kehilangan sebagian ketebalannya karena pemadatan.

\textbf{18.} Jumlahnya 1544, sehingga rerata aritmetiknya adalah \ang{128.7}, arah tenggara: tidak bermakna, karena bacaan tersebar di sekitar utara dan nilai dekat 340 serta dekat 010 dirata-ratakan seolah berjauhan.

\textbf{19.} $S = \sum\sin\alpha_i = 1.698$ dan $C = \sum\cos\alpha_i =
11.248$.

\textbf{20.} $\bar R = \sqrt{1.698^2 + 11.248^2}/12 = 11.376/12 = 0.948$: bacaan-bacaan itu berkelompok rapat di sekitar satu arah.

\textbf{21.} Sungai-sungai mengalir, dan delta membangun, ke arah sekitar utara (\cref{prop:b1:clastic-sediments:azimuths}); daratan berada di selatan.

\textbf{22.} Karena foreset diukur pada lapisan yang tidak lagi mendatar: \omterm{def:b1:geological-maps:attitude}{arah kemiringan} mereka adalah arah lapisan yang miring ditambah \omterm{def:b1:geological-maps:attitude}{kemiringan} itu sendiri. Memutar setiap foreset terhadap \omterm{def:b1:geological-maps:attitude}{jurus} perlapisan sebesar \omterm{def:b1:geological-maps:attitude}{kemiringan} perlapisan memulihkan perlapisan, dan bersamanya foreset, ke posisi asalnya.

\textbf{23.} Arus pasang surut mengalir masuk dan keluar, dan membangun set yang miring ke dua arah yang berlawanan, seperti dalam \omterm{def:b1:clastic-sediments:cross-bedding}{perlapisan silang} tulang ikan; diagram mawarnya akan memiliki dua lobus yang berjarak sekitar \ang{180}. Dua belas bacaan itu membentuk satu lobus, $\bar R$ mendekati 1: arus sungai yang mengalir ke satu arah.

\textbf{24.} $\bar\alpha = \operatorname{atan2}(1.698, 11.248) = \ang{8.6}$: \omterm{def:b1:spherical-astronomy:horizontal}{azimut} \omterm{def:b1:clastic-sediments:cross-bedding}{paleoarus} rerata satuan~3 adalah \ang{9}, sedikit ke timur dari utara.
\end{solution}
